\section*{ID7771: Half-Step Bessel Drop Operator \textbackslash{}(\textbackslash{}Delta\{\textbackslash{}text\{Bessel\}\}\textbackslash{}): dz\textasciicircum{}ν}
\paragraph{Metadata.}
PiID: 0299607838690 | RTEID: RTE:P30176.5188:T5.7393 | UUID: 7659a4bc-b358-5509-8738-63dc5b489403

\paragraph{Canonical Statement.}
Numerically, we can implement a truncated sum and represent \textbackslash{}(\textbackslash{}mathrm\{d\}z\textasciicircum{}\textbackslash{}nu\textbackslash{}) as finite difference for the appropriate power/derivative weight. Practically one will compute the spectral coefficients \textbackslash{}(c\_\textbackslash{}nu = \textbackslash{}langle f, \textbackslash{}psi\_\textbackslash{}nu\textbackslash{}rangle\textbackslash{}) with basis \textbackslash{}(\textbackslash{}psi\_\textbackslash{}nu(r) = J\_\{\textbackslash{}nu\}(k\_\textbackslash{}nu r)\textbackslash{}), then map \textbackslash{}(c\_\textbackslash{}nu \textbackslash{}mapsto (J\_\{\textbackslash{}nu+1/2\}-J\_\textbackslash{}nu) c\_\textbackslash{}nu\textbackslash{}).

\paragraph{Canonical Equation.}
\[
\mathrm{d}z^\nu
\]

\paragraph{Dictionary.}
Numerically, we can implement a truncated sum and represent \textbackslash{}(dz\textasciicircum{} \textbackslash{}) as finite difference for the appropriate power/derivative weight.

\paragraph{Encyclopedia.}
Half-Step Bessel Drop Operator \textbackslash{}(\textbackslash{}Delta\{\textbackslash{}text\{Bessel\}\}\textbackslash{}): dz\textasciicircum{}ν is a symbolic expression, written dz\textasciicircum{}ν. It introduces a symbol or relation used elsewhere in the framework. It is built from ν, z. Within the Trawin Zero Point registry it serves as one element of the framework's mathematical narrative.
